\section{Introduction}\label{sec:introduction}

Fix a prime \(p\).  Write \(L_i=L_{E(i)}\) for localization at
Johnson--Wilson theory and \(L_{K(i)}\) for localization at Morava
\(K\)-theory.  We use the derived \(p\)-completion
\(S=\mathbb S_p^\wedge\) of the sphere and study the overlap
\[
 T=L_{K(3)}S,\qquad X=L_2T.
\]
The object \(X\) is the part of the height-three local sphere visible
below height three.  Our question is how to build it from lower local
spheres while retaining the maps by which the pieces are joined.

Chromatic fracture makes the gluing problem precise.  An
\(E(2)\)-local spectrum is a homotopy pullback of its \(K(2)\)-local
and \(E(1)\)-local parts over their common localization.  The gluing
map is part of the data; the fracture square does not determine it
from the two corners alone
\cite[Section~2.4, diagram~(2.19)]{BB}.
Hopkins' chromatic splitting conjecture, as recorded by Hovey
\cite[Conjecture~4.2]{Hovey}, proposes classes, factorizations, and
summands that would split the relevant canonical cofiber sequence.
The revised strong formulation of Barthel--Beaudry
\cite[Conjecture~6.3 and Remark~6.4]{BB} has the same lower-sphere
pattern at the height and primes considered here.  We study an ordered
filtration with those cofibers and a canonical first unit, allowing
nonzero connecting maps.

The low-height results already distinguish these formulations.
Hovey's July 1993 account records the height-one case and the
height-two case at primes greater than three, attributing the latter
to Hopkins using Shimomura--Yabe \cite[pp.~17--19]{Hovey}.
Goerss--Henn--Mahowald later proved
the height-two splitting at the prime three
\cite[Theorem~1.2]{GoerssHennMahowald2014}.  Their proof identifies
the actual rational comparison map before applying chromatic fracture
\cite[proof of Theorem~5.11, p.~1288]{GoerssHennMahowald2014}.
At height two and the prime two, Beaudry disproved the original strong
form \cite[Theorem~1.4]{Beaudry2017}; Beaudry--Goerss--Henn established
a corrected decomposition with Moore-spectrum terms that retains the
weak unit splitting \cite[Theorem~1.1.6]{BGH2022}.

Two established calculations explain the pieces at height three.
At odd primes the height-two splitting gives two \(E(1)\)-local
pieces and two rational pieces
\cite[Theorem~6.7]{BB}.  At every prime and positive height, the theorem of
Barthel--Schlank--Stapleton--Weinstein computes the rational homotopy
of the \(K(n)\)-local sphere as an exterior algebra over \(\Qp\)
on generators of degrees \(1-2i\), \(1\leq i\leq n\)
\cite[Theorem~A]{BSSW}.  At height three the resulting degrees are
\[
 0,\ -1,\ -3,\ -4,\ -5,\ -6,\ -8,\ -9.
\]
This algebra calculation determines neither integral representatives
nor their images under lower-height localization.  Those map
identifications are what allow the layers to be assembled.

\subsection{The explicit filtration}

All spectra and maps in the theorem carry their natural \(S\)-module
structures.  A filtration means a sequence of homotopy cofiber
sequences: its \(i\)-th cofiber \(C_i\) has a connecting map
\(\partial_i:C_i\to\Sigma F_{i-1}\).

\begin{theorem}[Explicit height-three filtration]\label{thm:main}
For every prime \(p\geq5\), there are \(E(2)\)-local \(S\)-modules
and maps
\[
 0=F_0\longrightarrow F_1\longrightarrow\cdots\longrightarrow F_8,
 \qquad e:F_8\xrightarrow{\simeq}X,
\]
whose successive cofibers \(C_i=\cofib(F_{i-1}\to F_i)\), in order,
are
\[
\begin{array}{c|l@{\qquad\qquad}c|l}
i&C_i&i&C_i\\ \hline
1&L_2S&5&\Sigma^{-5}H\Qp\\
2&\Sigma^{-1}L_2S&6&\Sigma^{-6}H\Qp\\
3&\Sigma^{-3}L_1S&7&\Sigma^{-8}H\Qp\\
4&\Sigma^{-4}L_1S&8&\Sigma^{-9}H\Qp.
\end{array}
\]
Under \(F_1\simeq L_2S\), the composite
\(F_1\to F_8\xrightarrow e X\) is the canonical map
\(L_2(S\to L_{K(3)}S)\).
\end{theorem}

The attachment identification is a separate formal part of the
result.  Once the stages and their fracture data have been defined,
\Cref{thm:attachment-identification} collects every connecting map
and both signed roofs.  The application in \Cref{sec:application}
establishes \Cref{thm:main} and the height-three instance of
\Cref{thm:attachment-identification} from one choice of local maps,
compatibility homotopy, and coherent inverse data.  The first unit is
canonical; the attachment identification is relative to those choices.

Existence of a canonical-unit filtration with these cofibers is also
the height-three case of
\cite[Theorem~1.1]{CompanionFilteredChromaticSplitting}, whose range
\(p>n+1\) becomes \(p\geq5\) here.  That theorem supplies one family
of product bases at all lower positive heights.  The present proof
constructs the height-three maps independently of that theorem,
identifies the rational degree-three comparison, and records every
attachment relative to one compatible set of choices.  Neither
filtration statement asserts the strong wedge or splits its first unit.
The construction of \Cref{thm:main} is independent of the
rational-obstruction theorem used in \Cref{sec:attachment}.  That
separate theorem supplies additional information about the
extensions: for this explicit filtration, the first height-one
connecting map is nonzero.

\subsection{Three blocks and two gluing problems}

Put \(Q=L_0S\simeq H\Qp\), and group the layers as
\[
\begin{aligned}
 W&=L_2S\vee\Sigma^{-1}L_2S,\\
 U_{\mathrm{mid}}&=\Sigma^{-3}L_1S\vee\Sigma^{-4}L_1S,\\
 V&=\Sigma^{-5}Q\vee\Sigma^{-6}Q
       \vee\Sigma^{-8}Q\vee\Sigma^{-9}Q.
\end{aligned}
\]
The first localized basis is an actual map
\[
 w=(1,\zeta):W\longrightarrow X,
\]
where \(\zeta:\Sigma^{-1}S\to T\) is the integral determinant
class.  It becomes an equivalence after \(K(2)\)-localization.
Hence its cofiber \(Y\) is \(E(1)\)-local.  To insert the two
summands of \(U_{\mathrm{mid}}\), one needs a map to \(Y\), not merely the
knowledge that \(Y\) has the expected local homotopy groups.

There are two descriptions of that map.  At height one, an integral
orientation class \(\delta\in\pi_{-3}L_{K(1)}T\) participates in
the actual basis
\[
 (1,\zeta,\delta,\zeta\delta):
 R_1\vee\Sigma^{-1}R_1\vee\Sigma^{-3}R_1\vee\Sigma^{-4}R_1
 \xrightarrow{\simeq}L_{K(1)}T,\qquad R_1=L_{K(1)}S.
\]
The first two components are localizations of the same maps used
in \(w\); quotienting them gives the \(K(1)\)-local part of a map
from \(U_{\mathrm{mid}}\).  Rationally, the degree-three primitive
\(\alpha_3\in\pi_{-3}L_0T\) maps to \(c\delta\) for a nonzero
\(c\in\Qp\), and its product with \(\zeta\) maps to
\(c\zeta\delta\).  Rescaling only the rational primitive by
\(c^{-1}\) makes the two descriptions agree on the rationalized
\(K(1)\)-local overlap.  The height-one fracture square then
constructs
\[
 h:U_{\mathrm{mid}}\longrightarrow Y.
\]
The specified compatibility homotopy, rather than a dimension count,
is the input to this gluing step.

Pulling \(X\to Y\) back along the first summand and then all of \(U_{\mathrm{mid}}\)
gives \(F_3\) and \(F_4\).  The cofiber of \(F_4\to X\) is rational.
The four remaining exterior products give an actual equivalence
from \(V\) to that cofiber, and their ordered partial sums give
\(F_5,\ldots,F_8\) by pullback.  The terminal projection is an
equivalence because the full map from \(V\) is.  The proof in
\Cref{sec:fracture} follows the two fracture squares far enough to
identify their connecting maps, including both rotation signs.

\subsection{Why the localized bases require geometry}

The two coefficient comparisons use the loci of the height-three
deformation where the connected part has height two and height one.
Their \'etale parts have heights one and two, respectively; these are
the coheight-one and ordinary strata.  Finite marking towers on these
loci are related to extensions of vector bundles on the
Fargues--Fontaine curve.  The bundle and local-system framework
comes from Fargues--Fontaine and Scholze--Weinstein
\cite{FarguesFontaine,ScholzeWeinstein}; relative full faithfulness
of derived sections is supplied by Ansch\"utz--Le Bras
\cite[Corollary~3.11]{AnschuetzLeBrasFourier}.
Retaining the integral Tate frames and the determinant lattice makes
the group actions and the determinant product identity part of the
comparison.  These data are later used to identify the coefficient
classes of the unit and determinant maps in \(w=(1,\zeta)\).

The Kummer and translation calculations in the coheight-one work of
Barthel--Mann--Ray--Schlank--Senger--Weinstein--Zhou
\cite[Sections~3.3--3.6]{BMR} provide methodological context for
the cohomology of the completed strata.  The geometric answers must
then be carried to the algebraic coefficients in ordinary Morava
cooperations.
A cochain in their marking union must occur at one finite stage on
its compact domain; a cochain with Laurent coefficients must have
one pole bound.  Completion does not preserve those requirements by
definition, so the comparisons must establish the required finite
bounds before returning to the algebraic coefficients.

On the locus with connected height two, the completed calculation
controls cohomology at each finite stage of the connected marking tower.
The comparison in \Cref{h3:sec:coheight-one} proves finite cohomology
for a cofinal family of finite products of complete local fields, then
removes completion while preserving the group actions and the constants
map.  The descent calculation in \Cref{h3:sec:descent} then uses these
algebraic coefficients to prove that the actual unit
and determinant map \(w=(1,\zeta)\) is a \(K(2)\)-equivalence.

On the locus with connected height one, the completed calculation
provides the constants comparison determined by the Tate frame of rank
two and the connected part.
To reach the ordinary coefficients, \Cref{h3:sec:ordinary} first proves
finite cohomology for each compact lattice of power series.  An
intermediate chromatic descent calculation then supplies finite cohomology
for the union of coefficients with bounded poles.  Only after this step is
the analytic completion removed.  The resulting constants comparison
is the input to \Cref{h3:sec:integral-basis}, where finite Witt
coefficients carry the orientation through compatible length transitions
and derived completion produces \(\delta\).

The rational comparison uses these integral maps and the established
rational exterior algebra, whose height-three proof is reconstructed
in \Cref{h3:app:rational-boundary} following \cite{BSSW}.
\Cref{h3:sec:rational} transports a trace class in degree three along
the actual Tate frame of rank two on the ordinary stratum.  The finite
Witt comparisons carry this class through the coefficient tower, whose
derived limit is taken before rationalization.  The resulting comparison
proves that its image in the line \(\Qp\delta\) is nonzero.  This supplies
the scalar \(c\) in the gluing overview.
Product naturality for the same \(\zeta\) gives the same scalar on
\(\zeta\alpha_3\).

\subsection{Reading the proof}

\Cref{sec:fracture} first constructs the ordered filtration from
precisely stated local maps and a compatibility homotopy, and identifies
all its attachments for those same choices.  The remainder of the paper
constructs the required maps.

\Cref{h3:sec:framework} fixes the coefficient, group, and cochain
conventions.  \Cref{h3:sec:towers} constructs the finite Tate-coordinate
towers and their integral frame comparisons, and
\Cref{h3:sec:analytic} computes the completed strata with their actual
constants maps.  \Cref{h3:sec:coheight-one} removes completion on the
coheight-one stratum, using finite local-field stages and the full
reduced-norm quotient.  \Cref{h3:sec:descent} then compares ordinary
residue cooperations with those stages and obtains the actual
height-two basis.

\Cref{h3:sec:ordinary} proves the second coefficient comparison in the
order required by its dependencies: root and Laurent-tail control,
compact lattice finiteness, intermediate chromatic finiteness,
bounded-pole finiteness, and only then the natural ordinary comparison.
\Cref{h3:sec:integral-basis} carries the orientation through finite Witt
coefficients and derived completion to the actual height-one basis.
The trace argument of \Cref{h3:sec:rational} identifies the localization of the
particular degree-three primitive.

Finally, \Cref{sec:application} inserts these maps into the earlier
fracture construction, proving the main theorem and its attachment
identification from one common choice.  \Cref{sec:attachment} proves a
conditional criterion for the first middle attachment and applies the
companion canonical-map theorem through the local completion,
localization, and scalar comparisons proved there.
\Cref{h3:app:rational-boundary} then reconstructs the rational exterior
algebra from its arithmetic and geometric inputs.
